\documentclass[12pt]{article}
\usepackage[utf8]{inputenc}
\usepackage[T1]{fontenc}
\usepackage{hyperref}
\title{Tippera Language Resources for K-12 Curriculum Development}
\author{Compiled by Claude Sonnet 5 (claude-sonnet-5) for Andr\'as Kornai}
\date{}
\begin{document}
\maketitle

\section{Introduction}
Tippera (ISO 639-3 tpe) is a Sino-Tibetan (Tibeto-Burman) language, also called Tipra, Tripuri or Kokborok. The search found almost nothing about Bangladesh, and nearly all results concern Kokborok in Tripura, India. They are listed here as related context, not as Bangladeshi resources, and the match between those materials and the tpe variety was not verified.

\section{Primary Educational Resources}
\begin{itemize}
\item Tripura government document on Kokborok in schools (one source cited by the search summary reports 45 primary, 49 secondary and 22 higher secondary schools): \url{https://ica.tripura.gov.in/sites/default/files/9923_22.pdf}
\item News report quoting the Tripura Chief Minister on Kokborok teaching (1,296 primary, 115 secondary and 65 higher secondary schools): \url{https://nenow.in/north-east-news/tripura/govt-committed-to-preserving-minority-languages-tripura-cm.html}
\end{itemize}

\section{Linguistic and Reference Materials}
\begin{itemize}
\item Wikipedia, Kokborok: \url{https://en.wikipedia.org/wiki/Kokborok}
\end{itemize}
No dictionary or grammar was found in search results.

\section{Religious and Community Resources}
None found in search.

\section{Digital and Audio Resources}
None found in search.

\section{Curriculum Development Context}
The main issue in Tripura is script. Kokborok has no traditional script of its own; students write the language paper in Bengali script, and Roman script has been pushed by Tipra Motha. Reports say the board now allows Roman script in examinations.
\begin{itemize}
\item \url{https://compass.rauias.com/current-affairs/tripura-board-allows-kokborok-roman-script-examinations/}
\item \url{https://www.theweek.in/wire-updates/national/2024/02/09/ces6-tr-kokborok-ld-debbarma.html}
\end{itemize}
Bangladesh-specific mother-tongue education material for Tippera was not found and needs a separate search.
\end{document}